\section{Method Development Results}
\label{sec:development_results}

\subsection{Adaptive Objective and Compute}

All four planner arms completed all 60 development evaluations, and no native
sample exceeded 15~N. Table~\ref{tab:component_ablation_dev} summarizes the
matched results across the three force targets. M4 has the lowest mean MRE and
NRMSE at every target. Objective adaptation and compute adaptation each
recover part of the difference from M1, while their combination gives the
largest 12~N improvement. M4 reduces 12~N RMSE from 2.927 to 2.600~N and raises
mean force from 9.308 to 9.714~N. The strict 12~N tracking gate nevertheless
remains 0/5 for every arm.

\begin{table*}[t]
\centering
\caption{Matched planner-component development ablation. MRE and NRMSE are
averaged over 15 seed--target evaluations per arm.}
\label{tab:component_ablation_dev}
\small
\begin{tabular}{lrrrr}
\toprule
Arm & MRE (\%) & NRMSE (\%) & Samples & Median time (ms) \\
\midrule
M1: fixed objective, fixed compute & 12.03 & 14.66 & 128.0 & 112.18 \\
M2: adaptive objective             & 10.94 & 13.77 & 128.0 & 107.86 \\
M3: adaptive compute               & 10.26 & 13.52 & 121.2 &  80.58 \\
M4: complete method                &  9.65 & 12.94 & 120.4 &  83.25 \\
\bottomrule
\end{tabular}
\end{table*}

The compute schedule lowers the mean sample count because stable acquisition
and lower-force tracking often use a smaller budget. At 12~N, the schedule
increases the sample count when predicted transient risk rises. The complete
method therefore allocates computation by regime rather than reducing it
uniformly.

\subsection{Behaviour-Cloning Strength}

The one-step force-prediction RMSE is effectively unchanged across BC
coefficients, whereas actor--teacher action RMSE falls from 1.006 without BC
to 0.099 and 0.090 at coefficients 0.5 and 2. Table~\ref{tab:bc_ablation_dev}
shows the resulting closed-loop difference. Without BC, none of the 15 runs
establishes contact. A coefficient of 0.5 restores contact but remains
seed-sensitive, with three task failures and two samples above 15~N. The
coefficient of 2 completes all runs without a sampled violation. BC therefore
stabilizes the actor proposal used by MPPI rather than improving the one-step
force model.

\begin{table*}[t]
\centering
\caption{Behaviour-cloning development ablation over five seeds and three
force targets.}
\label{tab:bc_ablation_dev}
\small
\begin{tabular}{rrrrrr}
\toprule
BC & Task & Tracking & No contact & Samples $>15$ & Max peak (N) \\
\midrule
0   &  0/15 &  0/15 & 15/15 & 0 &  0.000 \\
0.5 & 12/15 & 10/15 &  0/15 & 2 & 16.513 \\
2   & 15/15 & 10/15 &  0/15 & 0 & 13.924 \\
\bottomrule
\end{tabular}
\end{table*}

\subsection{Router and Risk-Representation Ablations}

The historical hard router was evaluated at six activation fractions from
0.70 to 0.95 across the same five checkpoints and three force targets. All
90 runs completed without a sampled force violation. The actor-step fraction
rose from 27.1\% to 60.5\%, whereas mean MRE changed gradually from 6.45\% to
5.25\%. At 12~N, actor use rose from 25.2\% to 88.1\%, yet mean MRE remained
10.16\% at the highest threshold. There is therefore no discrete operating
point that resolves the high-force limitation; the threshold principally
controls planner use.

The earlier registered-frame controller also permits a sample-pooled
decomposition of its 56,481 contact samples at 12~N. Stable MPPI operation
accounts for 51,771 samples with a mean force of 11.020~N, actor re-entry
accounts for 1,950 samples at 10.075~N, and initial or other contact accounts
for 2,760 samples at 6.605~N. Their pooled mean is 10.772~N; the corresponding
mean of per-evaluation contact means is 10.767~N. The stable planner segment
therefore contains most of the negative bias, while actor re-entry deepens it.
This evidence motivated separate router and risk-representation ablations
rather than attributing the 12~N shortfall to contact loss alone.

The three risk representations each complete 15/15 runs, pass strict tracking
in 10/15, and produce no sample above 15~N. Force-only and point-envelope
trajectories are identical. The calibrated alternative changes only the
12~N trajectories and gives a slightly higher aggregate NRMSE (12.94\% versus
12.91\%). The frozen rule therefore selects next-force prediction alone.
These development results select the full M4 objective-and-compute mechanism,
BC coefficient 2, and the force-only risk representation before the
factor-separated evaluation.
